\begin{table}[htbp]
\centering
\begin{threeparttable}
\caption{Profile moderators of attacked effectivity: controlled contrasts, SEM paths, and descriptive susceptibility weights.}
\label{tab:multivariate}
\small
\setlength{\tabcolsep}{4.5pt}
\renewcommand{\arraystretch}{1.10}
\begin{adjustbox}{max width=\linewidth}
\begin{tabular}{p{0.34\linewidth}rrrrrrrrr}
\toprule
Moderator & Role & Controlled mean b & Controlled p & Ridge mean b & Weight \% & SEM mean b & SEM mean |b| & SEM min p & SEM paths \\
\midrule
Big Five Openness To Experience Mean \% & core & 2.839 & 0.016 &  &  & 2.822 & 2.822 & 0.002 & 8.000 \\
Big Five Neuroticism Mean \% & core & 2.596 & 0.172 &  &  & 2.787 & 2.787 & 0.006 & 8.000 \\
Sex Female & core & 3.857 & 0.270 &  &  & 3.267 & 3.630 & 0.013 & 8.000 \\
Big Five Conscientiousness Mean \% & core & -1.191 & 0.522 &  &  & -1.472 & 1.726 & 0.016 & 8.000 \\
Age Years & core & 1.243 & 0.154 &  &  & 1.019 & 1.019 & 0.178 & 8.000 \\
Big Five Agreeableness Mean \% & exploratory & 1.514 & 0.287 &  &  &  &  &  &  \\
Big Five Extraversion Mean \% & exploratory & 0.339 & 0.801 &  &  &  &  &  &  \\
\bottomrule
\end{tabular}
\end{adjustbox}
\begin{tablenotes}[flushleft]
\footnotesize
\item Note. Controlled coefficients summarize moderator associations with mean attacked effectivity. SEM columns summarize the repeated-outcome path model across the attacked opinion leaves. Weight percentages come from the target-conditional regularized aggregation used to compute the post hoc susceptibility index.
\end{tablenotes}
\end{threeparttable}
\end{table}